\subsection{Matched External Framework Track and Claim Scope}
\label{sec:external}

The repository also contains a preregistered fixed-topology track that makes the
external-framework comparison explicit. It uses the same task text, PDK/model
profile, simulator, independent verifier, 300 simulation-unit cap, 12-call cap,
and seeds 0--2 for every row. The seven paper-backed rows are AnalogCoder
\cite{lai2025analogcoder}, ADO-LLM \cite{yin2024adollm}, AnaFlow
\cite{ahmadzadeh2025anaflow}, Artisan \cite{chen2024artisan}, AutoSizer
\cite{yu2026autosizer}, ATLAS \cite{oh2026atlas}, and GLOVA
\cite{kim2025glova}. The ATLAS row is a manifest-backed clean-room adapter
for the public ATLAS repository. It is not presented as an official
reproduction because the current checkout does not bundle the original paper
experiment artifacts or claim equivalence to the original GF180 setup.
The complete task, framework identity, and implementation-status manifests are
in \path{data/bench/paper_sizing_tasks.jsonl} and
\path{data/bench/paper_sizing_manifest.json}.

This track is broader than the legacy E1g table above and is reported as a
separate repository-level artifact. Its primary aggregate, raw shards, and
strict merger are \path{results/paper_sizing_multiseed.json},
\path{results/paper_sizing_multiseed.md}, and
\path{scripts/audit_xsota_gate.py}. The frozen SIZ7g SPECTRA row does not
strictly beat the matched AnaFlow row on all four preregistered endpoints, so
the literal XSOTA-1 gate remains not met. We retain the tie and negative
endpoint cells rather than selecting a favorable subset.

The public AnalogGym staged policy is a separate sensitivity arm. It is
pre-registered as 50 nominal warm-start evaluations followed by 50 candidates
screened on five fixed corners, with an independent 45-condition verifier.
The complete HoiLee and Leung three-seed receipts are
\path{results/analoggym_pvt_staged_comparison.md} and
\path{results/analoggym_leung_pvt_staged_comparison.md}. Those results
have higher point estimates for PVT success and lower normalized violation on the two declared topologies,
while TT success ties and wall time is higher. They are therefore sensitivity
evidence, not a replacement for the frozen primary row and not a universal
analog-sizing SOTA claim.

Public cross-PDK and cross-simulator tracks are reported only when their source
contract supplies comparable metrics. ATLAS/GF180, SKY130, ASAP7, PTM and BJT
receipts are kept in the public-scope tables. AnalogPortBench, SGAnalog,
Analog Design Bench and LOADBench contribute source-native or capability
receipts where a common SPECTRA/AnaFlow transfer protocol is not available.
Those rows document coverage and reproducibility boundaries and are not silently
converted into ranking cells.

\paragraph{Cross-simulator sensitivity.}
The public Xyce-7.10 track uses the deterministic AC translator in
\path{scripts/xyce_ac_translator.py} and the same independent 45-condition
verifier. The matched robust HoiLee row gives SPECTRA TT/PVT success
$\XyceHoiTT{}/\XyceHoiPVT{}$ versus AnaFlow $\XyceHoiAnaTT{}/\XyceHoiAnaPVT{}$, with mean PVT violation
$\XyceHoiTTViol{}/\XyceHoiPVTViol{}$ versus $\XyceHoiAnaTTViol{}/\XyceHoiAnaPVTViol{}$. On Leung, SPECTRA and AnaFlow reach
$\XyceLeungTT{}/\XyceLeungPVT{}$, while SPECTRA has a lower mean PVT violation point estimate
from $\XyceAnaLeungPVTViol{}$ to $\XyceLeungPVTViol{}$.
These are complete six-cell-per-topology sensitivity receipts, not a hidden
replacement for the frozen SIZ7g row. The HoiLee TT deficit and PVT-success
ties remain visible in the gate audit.

\paragraph{SGAnalog harness boundary.}
The native SGAnalog harness evaluates all \SGAnalogTasks{} pinned sizing tasks over three
GPT-6.1-sol seeds with the source testbenches. The separate Xyce AC probe
translates and attempts all \SGAnalogTasks{} decks, preserving source and translated hashes,
return codes and log digests in
\path{results/sganalog_xyce_ac_probe.{json,md}}. All \SGAnalogFailures{} attempts stop at the
same included mismatch-function expression that Xyce rejects, so this probe
records a reproducible simulator boundary and contributes no ranking metric.
The distinction matters because capability coverage alone cannot establish
cross-simulator sizing superiority.
